% Cluster-category positive-half gluing for Y^+.
% Primary sources:
%   Keller, B., Yang, D., "Derived equivalences from mutations of quivers
%     with potential", Adv. Math. 226 (2011) 2118-2168 (arXiv:0906.0761).
%   Iyama, O., Yoshino, Y., "Mutation in triangulated categories and
%     rigid Cohen-Macaulay modules", Invent. Math. 172 (2008) 117-168
%     (arXiv:math/0607736).
%   Amiot, C., "Cluster categories for algebras of global dimension 2
%     and quivers with potential", Ann. Inst. Fourier 59 (2009) 2525-2590
%     (arXiv:0805.1035).
%   Fomin, S., Zelevinsky, A., "Cluster algebras II: finite type
%     classification", Invent. Math. 154 (2003) 63-121 (arXiv:math/0208229).
%   Ginzburg, V., "Calabi-Yau algebras" (arXiv:math/0612139).
%   Kontsevich, M., Soibelman, Y., "Cohomological Hall algebra, exponential
%     Hodge structures and motivic DT invariants" (arXiv:1006.2706).
%   Nagao, K., "Donaldson-Thomas theory and cluster algebras",
%     Duke Math. J. 162 (2013) 1313-1367 (arXiv:1002.4884).

\providecommand{\Hom}{\mathrm{Hom}}
\providecommand{\End}{\mathrm{End}}
\providecommand{\Ext}{\mathrm{Ext}}
\providecommand{\per}{\mathrm{per}}
\providecommand{\CoHA}{\mathrm{CoHA}}
\providecommand{\Jac}{\mathrm{Jac}}
\providecommand{\Rep}{\mathrm{Rep}}
\providecommand{\mod}{\mathrm{mod}}
\providecommand{\fd}{\mathrm{fd}}
\providecommand{\sC}{\mathcal{C}}
\providecommand{\sD}{\mathcal{D}}
\providecommand{\sH}{\mathcal{H}}
\providecommand{\bC}{\mathbb{C}}
\providecommand{\bZ}{\mathbb{Z}}
\providecommand{\bP}{\mathbb{P}}
\providecommand{\Yplus}{Y^{+}}
\providecommand{\kch}{\kappa_{\mathrm{ch}}}
\providecommand{\kcat}{\kappa_{\mathrm{cat}}}
\providecommand{\chiHall}{\chi^{\mathrm{Hall}}}

\section*{Cluster Categories and Positive Hall Halves}
\label{sec:cluster-Yplus-gluing}

\paragraph{Scope.}
Let $(Q,W)$ be a reduced Jacobi-finite quiver with potential and let
$\Gamma(Q,W)$ be the complete Ginzburg Calabi-Yau $3$ dg algebra, so
\[
  H^0\Gamma(Q,W)=\Jac(Q,W).
\]
A cluster seed is the pair consisting of a Jacobi presentation
$(Q,W)$ and a stability chamber $\sigma$ whose heart is
$\sH_\sigma=\mod\Jac(Q,W)$.  An admissible mutation edge is a vertex
$k$ for which Derksen-Weyman-Zelevinsky mutation is defined, the reduced
mutated potential $(Q',W')=\mu_k(Q,W)$ is Jacobi-finite, the adjacent
heart is the simple tilt of $\sH_\sigma$ at $S_k$, and orientation data
is transported through the Keller-Yang equivalence.

The output of this construction is the transition map
\[
  g_k:
  \widehat{\Yplus}_{\sigma}(Q,W)
  \longrightarrow
  \widehat{\Yplus}_{\sigma'}(Q',W')
\]
between completed positive Hall halves.  Its mathematical content is
seed mutation plus Kontsevich-Soibelman wall crossing.  A full BKM
algebra, full Yangian, or $\mathcal W$-algebra enters after extra data:
a negative half, Cartan factor, non-degenerate Hall pairing, completion,
Drinfeld double or centre, and a representation or evaluation functor.
For Calabi-Yau dimension $d\geq3$, the native $\Phi_d$ output is
$E_1$-chiral; the $E_2$ layer lives at the centre, double, or module
level.

\subsection*{The Amiot Category}

Amiot associates to the Jacobi-finite Ginzburg algebra the quotient
\[
  \sC(Q,W)=\per(\Gamma(Q,W))\big/\sD_{\fd}(\Gamma(Q,W)).
\]
This category is Hom-finite and $2$-Calabi-Yau, and the image
$\pi\Gamma(Q,W)$ is a cluster tilting object.  For acyclic Dynkin $Q$
and $W=0$, the same construction recovers the classical cluster
category of the corresponding Fomin-Zelevinsky finite type.

The cluster category records tilting presentations.  The Hall algebra
uses the additional critical presentation.  For a dimension vector
$\gamma\in\bZ_{\geq0}^{Q_0}$ and
$G_\gamma=\prod_{i\in Q_0}GL_{\gamma_i}$, the cohomological summand is
\[
  \CoHA_\gamma(Q,W)
  =
  H^{\bullet}_{G_\gamma,c}
  \bigl(\Rep_\gamma(Q),
  \phi_{\mathrm{Tr}(W)_\gamma}\mathbb Q\bigr)
  \otimes \mathbb L^{\chi_Q(\gamma,\gamma)/2}.
\]
The Hall product is the vanishing-cycle correspondence for short exact
sequences in the heart $\sH_\sigma$.  The square root of the Tate twist
is fixed by orientation data.

\subsection*{Cluster Charts}

Iyama-Yoshino mutation applies to cluster tilting objects in a
$2$-Calabi-Yau triangulated category.  Inside the mutation class of
$\pi\Gamma(Q,W)$, a cluster tilting object $T$ gives
\[
  A_T=\End_{\sC(Q,W)}(T),
\]
and the Keller-Yang comparison identifies $A_T$ with the Jacobi algebra
of the corresponding mutated quiver with potential.  Thus a quiver
chart is the data
\[
  U=(T_U,Q_U,W_U,\sigma_U,o_U),
\]
where $o_U$ is orientation data and $\sigma_U$ is the chamber whose
heart is $\mod\Jac(Q_U,W_U)$.

The mutation groupoid has one elementary generator $\mu_k$ for each
mutable vertex and the usual exchange relations.  A path in this
groupoid records a finite sequence of adjacent chamber changes.

\subsection*{Keller-Yang Mutation}

For an admissible mutation at $k$, Keller and Yang construct triangle
equivalences
\[
  \mu_k^{+}:\per(\Gamma(Q,W))\xrightarrow{\sim}\per(\Gamma(Q',W')),
  \qquad
  \mu_k^{-}:\per(\Gamma(Q',W'))\xrightarrow{\sim}\per(\Gamma(Q,W)),
\]
restricting to
\[
  \sD_{\fd}(\Gamma(Q,W))\simeq\sD_{\fd}(\Gamma(Q',W')).
\]
They descend to an equivalence
\[
  \sC(Q,W)\simeq\sC(Q',W').
\]
Keller-Yang Theorem 6.3 identifies the endomorphism dg algebra of the
mutated tilting object with $\Gamma(Q',W')$.  On hearts, the operation
is the simple tilt at $S_k$; on charge lattices, it is the cluster
mutation of the basis of simples.

\subsection*{Positive Hall Transport}

Fix the positive cone $C_\sigma^+\subset\bZ^{Q_0}$ determined by
$\sH_\sigma$.  The completed positive half is
\[
  \widehat{\Yplus}_{\sigma}(Q,W)
  =
  \widehat{\bigoplus_{\gamma\in C_\sigma^+}
  \CoHA_{\gamma}(Q,W)}.
\]
The completion is taken in the positive charge cone and in the
motivic or cohomological grading used by the Hall product.

Kontsevich-Soibelman wall crossing compares two ordered factorizations
of the same motivic Hall element.  Nagao identifies this comparison
with cluster mutation for the chamber change above.  After choosing
orientation data and a cohomological realization, the wall element
gives an isomorphism of completed chamber algebras
\[
  g_k:
  \widehat{\Yplus}_{\sigma}(Q,W)
  \xrightarrow{\sim}
  \widehat{\Yplus}_{\sigma'}(Q',W').
\]
Equivalently, after the Kontsevich-Soibelman integration map, $g_k$ is
the cluster transformation of the completed quantum torus.

For a path $p:U\to V$ in the mutation groupoid, define
\[
  g_{UV}=g_p:
  \widehat{\Yplus}_{\sigma_U}(Q_U,W_U)
  \longrightarrow
  \widehat{\Yplus}_{\sigma_V}(Q_V,W_V).
\]
For fixed paths and orientation data, the transition maps satisfy
\[
  g_{UW}=g_{VW}\circ g_{UV}
\]
on triple overlaps.  A loop in the mutation groupoid acts by the
Picard-Lefschetz monodromy of the corresponding stability loop.  When
that monodromy is represented by a Hall wall element, the action is
conjugation by that element in the completed Hall algebra.

\subsection*{The Toric Normal Form}

For $\bC^3$, the critical quiver is the one-vertex three-loop quiver
with cubic potential
\[
  W_{\bC^3}=x[y,z].
\]
Kontsevich-Soibelman and Schiffmann-Vasserot identify its cohomological
Hall algebra with the positive half of the affine Yangian:
\[
  \CoHA(\bC^3)\cong\Yplus(\widehat{\mathfrak{gl}}_1).
\]
This is an associative $E_1$ positive half.  The full affine Yangian
appears after adjoining the negative half and Cartan factor through
the completed Hall-Drinfeld double:
\[
  \Yplus(\widehat{\mathfrak{gl}}_1)
  \hookrightarrow
  \widehat D_{\mathrm{Hall}}
  \bigl(\Yplus(\widehat{\mathfrak{gl}}_1)\bigr)
  =
  Y^-\widehat\otimes Y^0\widehat\otimes \Yplus.
\]
A Fock representation or vacuum evaluation is a further step:
\[
  \widehat D_{\mathrm{Hall}}
  \bigl(\Yplus(\widehat{\mathfrak{gl}}_1)\bigr)
  \longrightarrow
  \End\bigl(\mathcal W_{1+\infty}[\lambda]\text{-vacuum}\bigr).
\]
The cluster atlas supplies the positive-half input for this chain.

\subsection*{Conifold}

The resolved conifold is
\[
  X_{\mathrm{con}}=
  \mathrm{Tot}\bigl(\mathcal O(-1)^{\oplus2}\to\bP^1\bigr).
\]
Its noncommutative crepant resolution is the Klebanov-Witten Jacobi
algebra: two vertices, arrows
$a_1,a_2:0\to1$, $b_1,b_2:1\to0$, and potential
\[
  W_{\mathrm{KW}}
  =
  a_1b_1a_2b_2-a_1b_2a_2b_1.
\]
The flop is the mutation at either vertex.  Nagao's theorem identifies
the induced wall crossing with the conifold DT wall-crossing operator.
The Hall object obtained here is the conifold positive half in the
chosen chamber; its double or centre is additional structure.

The chiral scalar attached to this Hall shadow is
\[
  \kch(A_{\mathrm{con}})=1.
\]
This value is the direct McKay/Jacobi-NCCR value of the conifold
chamber.  The conifold is a rank-two bundle over a curve, so the local
surface formula for $\mathrm{Tot}(\omega_S)$ is absent from this
calculation.

\subsection*{Local $\bP^2$}

Let
\[
  X_{\bP^2}=\mathrm{Tot}(\omega_{\bP^2}).
\]
The $\bZ/3$ McKay presentation has three vertices and the cyclic
Beilinson quiver with three arrows between consecutive vertices.  Its
Jacobi potential is cubic.  The three toric affine charts form the
affine $\widehat A_2$ mutation cycle, and the positive-half gluing is
the corresponding cluster-groupoid cocycle acting on
\[
  \widehat{\Yplus}_{\sigma}(X_{\bP^2}).
\]
The Stage-1 chiral scalar is
\[
  \kch(A_{X_{\bP^2}})=\frac{1}{2}\chi_{\mathrm{top}}(\bP^2)=\frac32.
\]
This is the compactly supported McKay shadow of the local surface
chart.

\subsection*{Typed Local Constants}

The local constants below are separated by construction:
\[
\begin{array}{c|c|c|c|c}
X & \text{positive Hall seed} & \kch & \kcat
  & \kappa_{\mathrm{BKM}}\\
\hline
\bC^3 &
\Yplus(\widehat{\mathfrak{gl}}_1) &
1 & 0 & \mathrm{n/a}\\
X_{\mathrm{con}} &
\text{Klebanov-Witten Jacobi CoHA} &
1 & 0 & \mathrm{n/a}\\
X_{\bP^2} &
\bZ/3\text{-McKay Jacobi CoHA} &
3/2 & 0 & \mathrm{n/a}
\end{array}
\]
The table records local Hall-half constants.  Borcherds products and
fibre lattice corrections live in separate constructions.  The BKM
invariant $\kappa_{\mathrm{BKM}}$ belongs to the Borcherds-product
construction, for example the Gritsenko-Clery rows, outside cluster
mutation of these positive Hall seeds.

\subsection*{Bridge To The Chiral Target}

For a fixed oriented cluster atlas, the construction above gives a
diagram in completed $E_1$ Hall algebras:
\[
  \operatorname*{hocolim}_{T\in\mathrm{Clust}(Q,W)}
  \widehat{\Yplus}_{\sigma_T}(Q_T,W_T).
\]
Identification with a chiral algebra produced by $\Phi_3$ requires the
following additional data:
\[
\begin{gathered}
  \text{finite-type critical charts with compactly supported or
  equivariant vanishing cycles,}\\
  \text{compatible orientation data and a Kontsevich-Soibelman
  integration map,}\\
  \text{degree-preserving critical-cohomology identifications along
  admissible mutations,}\\
  \text{Hall multiplication compatible with all transition maps,}\\
  \text{a negative half, Cartan factor, non-degenerate Hall pairing,
  and completed double,}\\
  \text{a centre or evaluation functor to the vertex-algebra target.}
\end{gathered}
\]
These inputs carry the positive Hall atlas to the chiral target.  The
CY$_3$ output itself remains $E_1$; the braided $E_2$ structure belongs
to the centre, double, or representation category.

\subsection*{Wall Crossing And Euler Series}

Define the completed Hall Euler series in the quantum torus of the
charge lattice by
\[
  \chiHall_{\sigma}(Q,W)
  =
  \sum_{\gamma\in C_{\sigma}^{+}}
  \left(
    \sum_i (-1)^i
    \dim H^i_{G_\gamma,c}
    \bigl(\Rep_\gamma(Q),
    \phi_{\mathrm{Tr}(W)_\gamma}\mathbb Q\bigr)
  \right)x^\gamma.
\]
For an admissible mutation edge, the Kontsevich-Soibelman
transformation gives
\[
  \chiHall_{\sigma'}(Q',W')
  =
  \mu_k^{\mathrm{KS}}\bigl(\chiHall_{\sigma}(Q,W)\bigr)
\]
after the cluster identification of charge lattices.  This transported
equality is the wall-crossing theorem.  Coefficientwise equality is a
separate fixed-sector statement in transported variables.

When the bridge data identify the transported Hall Euler scalar with
the modular characteristic of the $\Phi_3$ output, the same equality
gives mutation invariance of $\kch$ for the constructed chiral object.
For the local geometries above, the invariant values are exactly
\[
  \kch(A_{\bC^3})=1,\qquad
  \kch(A_{\mathrm{con}})=1,\qquad
  \kch(A_{X_{\bP^2}})=\frac32.
\]
